\subsection{The complete quartic multiplier audit}

\statement{C.20}{Proved}{The actual dihedral quotient behind the quartic argument}
Let $\alpha$ be the actual quartic torus representative of
C.19, $c=[X,Y]$, and $E=F\rtimes\langle s\rangle$.
There is a unique nonzero invariant character
$\chi:F\to C_2$.  One can choose a discrete free basis
$x,y$, with $[x,y]=c$, $\chi(x)=1$, $\chi(y)=0$, such
that the full quotient
\[
 F/\overline{\langle\!\langle x^2,(xy)^2\rangle\!\rangle}
   =\widehat{\mathbb Z}\rtimes C_2
   =\langle r,t\mid t^2=1,\ trt=r^{-1}\rangle,
 \qquad x\mapsto t,\quad y\mapsto r                         \tag{C.20.1}
\]
is $\alpha$-invariant.  The induced action is
\[
 \bar\alpha(r)=r^{-\lambda},\quad
 \bar\alpha(t)=r^jt,\quad (1-\lambda)j=0,
                         \quad j_2\in\mathbb Z_2^\times.   \tag{C.20.2}
\]
At every negative odd coordinate $q$, this gives a
specified quotient $\rho:F\to D_{2q}$ with
$\rho\alpha=\rho$, $\rho(c)=r^{-2}$, and
$q\equiv1\pmod4$.  This independently proves the
portion of LM4 D.68 needed for D.70.

\begin{proof}
By C.19, $A_2^2=-I$.  Its reduction cannot be identity:
$(I+2M)^2\equiv I\pmod4$, whereas $-I\not\equiv I$.
Thus it is a nonidentity transvection modulo two and
has a unique nonzero fixed dual character $\chi$.
For its possible values on $X,Y$, choose respectively
\[
\begin{array}{c|cc}
(\chi(X),\chi(Y))&x&y\\\hline
(1,0)&X&Y\\
(0,1)&XYX^{-1}&X^{-1}\\
(1,1)&X&YX^{-1}.
\end{array}
\]
The inverse substitutions in the last two rows are
$(X,Y)=(y^{-1},yxy^{-1})$ and $(x,yx)$.
Free reduction proves the stated commutator identity
in each case, so these are also inverse continuous
profinite basis maps.

Write $F_0=\ker\chi$.  Its type-$(1,2)$ covering
surface gives Schreier basis
$u=x^2$, $B=y$, $C=xyx^{-1}$ and cusp $c=CB^{-1}$.
On $M=H_1(F_0,\widehat{\mathbb Z})$, deck conjugation
by $x$ is the involution $J$ fixing $u$ and exchanging
$B,C$.  The inclusion in base homology is
\[
 I=\begin{pmatrix}2&0&0\\0&1&1\end{pmatrix},
 \qquad\ker I=\widehat{\mathbb Z}b,\quad b=C-B.
                                                               \tag{C.20.3}
\]
Choose a labelled cusp-normalizer element $v=f_0s$.
By applying C.19 to its square,
\[
 vcv^{-1}=c^\lambda,\qquad v^4=c^\kappa,
                           \quad\kappa_2\text{ odd}.        \tag{C.20.4}
\]
The character extends to $E$ by $\widetilde\chi(s)=0$.
Let $T$ be conjugation by $v$ on $M$.  Since
$\chi(vxv^{-1})=\chi(x)$, one has $TJ=JT$.
On the boundary line $T$ acts by $\lambda$; on
the filled handle quotient it acts by $A$.
Rationally split these $J$ eigenspaces at each prime.
The identities $A^2=-I$ and $\lambda^2=1$ give
$T^2=-J$, first rationally and hence on the torsion-free
integral module.  At two, $Tb=b$, so
\[
 1+T+T^2+T^3=2(1-J),\qquad
 (1+T+T^2+T^3)M_2=2\mathbb Z_2b.                           \tag{C.20.5}
\]
If $\widetilde\chi(v)=0$, write $v=ns$ with $n\in F_0$.
Then $(ns)^4=n\alpha(n)\alpha^2(n)\alpha^3(n)$,
whose homology belongs to the right side of (C.20.5).
This contradicts $[v^4]=\kappa_2b$.  Hence
$\chi(f_0)=1$.  If $S$ is the action of $\alpha$ on
$M$, then $T=JS$, $S^2=-J$, and $Sb=-\lambda b$.

For complete control of the character used next, write
$A=(a_{ij})$.  The equations $TJ=JT$, $Tb=\lambda b$
and $IT=AI$ give
\[
 T=\begin{pmatrix}
 a_{11}&a_{12}/2&a_{12}/2\\
 a_{21}&(a_{22}+\lambda)/2&(a_{22}-\lambda)/2\\
 a_{21}&(a_{22}-\lambda)/2&(a_{22}+\lambda)/2
 \end{pmatrix},\qquad S=JT.                               \tag{C.20.6}
\]
All entries are integral: $a_{12}$ is even and
$a_{22},\lambda$ are odd.  Define
$\ell:F_0\to\widehat{\mathbb Z}$ by
$\ell(u)=0$, $\ell(B)=1$, $\ell(C)=-1$.
Conjugation by $x$ changes its sign and (C.20.6) gives
$\ell\alpha=-\lambda\ell$.  Thus $\ker\ell$ is normal
in $F$ and $\alpha$-invariant.  Sending $x$ to $t$ and
$y$ to $r$ restricts to $n\mapsto r^{\ell(n)}$ on
$F_0$, so its exact kernel is $\ker\ell$.
The relations $x^2=(xy)^2=1$ conversely give
$xyx=y^{-1}$, proving both presentation maps in
(C.20.1).  This also verifies that its rotation
subgroup is full $\widehat{\mathbb Z}$.

The induced images have the first two forms in
(C.20.2).  Squaring gives $r\mapsto r$ and
$t\mapsto r^{(1-\lambda)j}t$; fourth power identity
and torsion-freeness of $\widehat{\mathbb Z}$ imply
$(1-\lambda)j=0$.  Write the image of $f_0$ as
$r^mt$.  The actual quotient of (C.20.4) yields
\[
 v^2=r^{(1+\lambda)m-j}s^2,\qquad
 v^4=r^{2((1+\lambda)m-j)}=r^{-2\kappa}.
\]
Thus $j=(1+\lambda)m+\kappa$, which is odd at two.
At a negative odd coordinate, $(1-\lambda)j=0$
gives $j_q=0$, and both generators of $D_{2q}$
are fixed.  Finally $A_q^2=-I$ and
$\det A_q=-1$ force $A_q$ to be a scalar primitive
fourth root: if the two roots of $X^2+1$ were
opposite, the determinant would be one.  Therefore
$\mathbb Q_q$ contains that fourth root, so
$q\equiv1\pmod4$.  This proves every asserted map
and scalar condition.
\end{proof}

\statement{C.21}{Proved}{The integral square on the nontrivial rotation summand}
Use a negative coordinate and the specified quotient of
C.20.  Put $N=\ker\rho$, $R=\langle r\rangle\cong C_q$,
$\beta=\alpha^2$, $\sigma=s^2$, $C=C_F(\sigma)$ and
$C_0=C\cap N$.  Then
\[
 \rho(C)=R,\qquad \dim H^1(C_0,\mathbb F_2)=q-1.
                                                               \tag{C.21.1}
\]
The cover $N$ has type $(q,2)$.  If
$V_2=H^1(N,\mathbb Z_2)$, $e=q^{-1}\sum_{r'\in R}r'$
and $W_2=(1-e)V_2$, then
\[
 \operatorname{rank}W_2=2(q-1),\qquad \beta^*|_{W_2}=I.
                                                               \tag{C.21.2}
\]
There are two section classes in $N\rtimes\langle\sigma\rangle$,
represented by $\langle\sigma\rangle$ and
$x\langle\sigma\rangle x^{-1}$, with normalizers
$C_0\times\langle\sigma\rangle$ and its $x$-conjugate.
This proves LM4 D.69, including its integral conclusion.

\begin{proof}
The exact quotient $E/N=D_{2q}\times C_4$ sends
$fs^j$ to $(\rho(f),s^j)$, since $\rho\alpha=\rho$.
The cusp monodromy $r^{-2}$ has two cycles of length
$q$ on its $2q$ sheets.  The cover has two cusps and
Euler characteristic $-2q$, hence genus $q$.
Its full completion is $N$ by finite-index subgroup
topology.  The intermediate type-$(1,2)$ group
$F_0=\rho^{-1}(R)$ has the Schreier basis of C.20,
with deck images $u\mapsto1$, $B\mapsto r$,
$C\mapsto r^{-1}$.  Put $\Lambda=\mathbb Z_2[R]$
and $\nu=\sum r'$.  The graph-cover cycle module is
the kernel of
\[
 \partial:\Lambda^3\to\Lambda,\qquad
          (a,b,d)\mapsto(r-1)b+(r^{-1}-1)d.
\]
It has explicit inverse parametrizations
\[
 (a,d,k)\mapsto(a,r^{-1}d+k\nu,d),\qquad
 (a,b,d)\mapsto(a,d,k),\quad b-r^{-1}d=k\nu.
\]
The equality defining $k$ is unique because
$\operatorname{ann}(r-1)=\mathbb Z_2\nu$, as follows
by comparing cyclic coefficients.  Thus homology,
and by duality its character module, is
$\mathbb Z_2\oplus\Lambda^2$.  This proves the rank
assertion in (C.21.2); restriction and odd-index
transfer identify $eV_2$ with $H^1(F_0,\mathbb Z_2)$.

Let $S$ be pullback by $\alpha^{-1}$ on $V_2$ and
$B=S^2=\beta^*$.  It commutes with all rotation deck
maps.  Over an unramified extension of $\mathbb Z_2$
containing the $q$th roots, every nontrivial character
summand is free of rank two.  Modulo two, $S^4=I$
makes $S$ unipotent, so its square is identity on
these rank-two summands.  Thus $B=I$ on
$W=W_2/2W_2$.  On $eV_2$ the dual of C.20 gives
$B=-J^{\mathsf T}$.  Its reduction has one Jordan
block of size two and one of size one, so on all
$V=H^1(N,\mathbb F_2)$ one has
\[
                       \operatorname{rank}(B-I)=1.         \tag{C.21.3}
\]
On the deck-trivial rank-three space the reduction
of $S$ has a single size-three block, since its
square is nonidentity.

The three full section classes of the positive
square have distinct labels in $H_1(F,\mathbb F_2)$,
by (C.19.1).  Use $\sigma$ as label zero.
Conjugation by $s$ acts on them by the linear
transvection $A_2\bmod2$, since it fixes $\sigma$.
An invariant three-element subset containing zero
in $\mathbb F_2^2$ omits the other fixed point of
this transvection.  Its other two points therefore
have $\chi$-value one.  Their sections cannot lie
in $N\rtimes\langle\sigma\rangle$.  Thus only the
full class of $\sigma$ occurs there.  Since
$H^1(C,\mathbb F_2)=0$, one has $\chi(C)=0$ and
$\rho(C)\leq R$, hence $\rho(C)=1$ or $R$.
The exact list of $N$-classes is the coset set
\[
\begin{split}
 D_{2q}/\rho(C)&\longrightarrow
 \{\text{section classes in }N\rtimes\langle\sigma\rangle\},\\
 \rho(f)\rho(C)&\longmapsto[f\langle\sigma\rangle f^{-1}].
\end{split}                                                   \tag{C.21.4}
\]
The normalizer criterion proves injectivity and
surjectivity, and conjugation by $s$ fixes this
coset set because $\rho\alpha=\rho$.  The same
set therefore gives the order-two subgroup classes
in $N\rtimes C_4$.

Suppose first that $\rho(C)=1$.  There are $2q$
classes in (C.21.4), each with normalizer conjugate
to $C\times C_2$.  Since $C$ has cohomological
2-dimension at most one and $H^1(C,\mathbb F_2)=0$,
Brown restriction identifies the localized
cohomology ring of $N\rtimes C_2$ with
$\mathcal R=\mathbb F_2^{2q}$.  All cohomology
in each degree is finite by the two-row spectral
sequence, so the dense restriction is onto.
For its quotient character $u$, every $\theta\in W$
extends by value zero on $\sigma$, and the split
two-row spectral sequence gives an injection
\[
                     W\longrightarrow\mathcal R,\qquad
                              \theta\longmapsto\theta/u.
                                                               \tag{C.21.5}
\]
Indeed section-induced injectivity on the bottom
row kills every possible differential, and the
nonzero row-one class remains nonzero after any
power of $u$.

Let $D$ be the reduction modulo two of $(I-B)/2$
on $W_2$.  For an integral lift
$\varphi:N\to\mathbb Z/4$ of $\theta$, the cochain
$\widetilde\theta(n\sigma^\epsilon)=\varphi(n)$ has
\[
 \delta\widetilde\theta(n\sigma^\epsilon,m\sigma^\eta)
       =\epsilon(\varphi(m)-\varphi(\beta(m)))
       =2\epsilon(D\theta)(m).
\]
Its Bockstein gives the exact identity
$\theta^2=uD\theta$.  All elements of $\mathcal R$
are idempotent, so (C.21.5) forces $D=I$ and hence
$B\equiv-I\pmod4$ on $W_2$.
If $S\bmod2$ on any nontrivial rank-two rotation
summand were identity, its integral square would
be $I\pmod4$, a contradiction.  Thus it has one
size-two Jordan block on each of the $q-1$
nontrivial summands, and one size-three block on
the trivial summand.  The periodic $C_4$ resolution
has maps $S-I$ and $(S-I)^3$; each such block
contributes one dimension in every positive degree.
Therefore $\dim H^j(C_4,V)=q$ for $j>0$, and
\[
 \dim H^n(N\rtimes C_4,\mathbb F_2)\leq q+1
                                      \quad(n\geq2).       \tag{C.21.6}
\]
But the $2q$ normalizers in that group are conjugates
of $C\rtimes C_4$, with cohomology equal to that
of $C_4$ by the vanishing just used.  Brown
restriction gives dimension $2q$ in high degrees,
contradicting (C.21.6).  Hence $\rho(C)=R$.

There are now two section classes, with the
normalizers claimed.  By (C.21.3), the periodic
$C_2$ cohomology of $V$ has dimension $2q-1$ in
positive degrees.  The split spectral sequence
therefore gives dimension $2q$ in high degrees
for $N\rtimes C_2$.  Brown restriction gives
$2q=2(1+\dim H^1(C_0,\mathbb F_2))$, proving
(C.21.1).  Its localized ring is
\[
\begin{split}
 \mathcal R\cong&
 (\mathbb F_2\oplus H^1(C_0,\mathbb F_2)u^{-1})\\
 &{}\times
 (\mathbb F_2\oplus H^1(xC_0x^{-1},\mathbb F_2)u^{-1}).
\end{split}
\]
Its nilradical $\mathcal I$ is square zero and
has dimension $2(q-1)$; the residue is
$\mathbb F_2^2$.  Choose $h\in C$ with $\rho(h)=r$.
Conjugation by $h$ fixes $\sigma$ and realizes
the rotation action on all the cohomology in
(C.21.5).  It fixes both residue factors, because
rotations fix the two cosets in $D_{2q}/R$.
Consequently the equivariant image of $W$ in
$\mathcal R/\mathcal I$ vanishes: averaging is zero
on $W$ and identity on the residue.  Thus (C.21.5)
identifies $W$ with the square-zero ideal
$\mathcal I$.  The same exact Bockstein identity
now gives $D=0$, or $B\equiv I\pmod4$ on $W_2$.
Write $B=I+4L$.  Since $B^2=I$, cancellation in
the torsion-free matrix algebra gives
$L(I+2L)=0$.  The second factor is invertible,
so $L=0$ and $B=I$.  This is the full integral
conclusion (C.21.2), not only its mod-two shadow.
\end{proof}

\statement{C.22}{Proved}{The quartic multiplier assertion of LM4 D.70}
Every order-four element of the full once-punctured-torus
loop-outer group has multiplier globally one.  Every
finite two-subgroup consequently has multiplier values
in the global units $\{1,-1\}$, determined by its
dyadic coordinate.  No geometric realization or
conjugacy assertion is included.

\begin{proof}
We first prove the fixed-character obstruction used
in the final comparison.  Let $P$ be the pro-two
completion of a closed orientable genus-$g$ surface,
$g\geq2$, and let $b$ be an actual involution which
is identity on $H^1(P,\mathbb F_2)$ but nonidentity
on integral dyadic homology.  In
$G=P\rtimes_b C_2$, every mod-two character of $P$
extends by value zero on the labelled section.
The quotient class $u$ survives, and the surface
class is a product of two extended characters.
These generate the second page
$\mathbb F_2[u]\otimes H^*(P,\mathbb F_2)$, so the
multiplicative spectral sequence collapses.  In
particular $u\eta$ survives in bidegree $(1,2)$
for a nonzero surface class $\eta$.

If $\theta,\psi:P\to\mathbb Z_2$ are integral
$b$-fixed characters, they extend to actual
$\mathbb Z_2$ characters of $G$, zero on the
labelled section.  They vanish on every finite
subgroup, by torsion-freeness of their target.
For every order-two section $T$, its action on
$P$ has the same nonidentity integral homology
as $b$.  Its fixed subgroup $C_P(T)$ is free
pro-two by Corollary~2 of
\cite{CEKochloukovaZalesskii2007}: $P$ is a
pro-two Poincare duality group of dimension two,
and this is an actual nontrivial order-two
automorphism, so the infinite-index step in that
corollary applies.  Its second cohomology is zero.
On $N_G(T)=C_P(T)\times T$, the two characters
are pulled back from $C_P(T)$; hence
$u\bar\theta\bar\psi$ restricts to zero there.
Brown restriction is injective in degree three
because the virtual dimension is two and there
is no $C_2^2$.  Thus $u\bar\theta\bar\psi=0$
in $H^3(G,\mathbb F_2)$.  Survival of $u\eta$
then implies
\[
 \bar\theta\smile\bar\psi=0
                 \quad\text{in }H^2(P,\mathbb F_2).        \tag{C.22.1}
\]
This is an assertion about reductions of integral
fixed characters; it does not say that all
mod-two fixed characters have zero pairing.

Take the actual quartic lift from C.19.  A negative
odd coordinate $q$ would give the quotient and
cover in C.20--C.21.  Fill its two cusps, represented
by $d_0=c^q$ and $d_1=xc^qx^{-1}$, using the exact
maps
\[
 N\twoheadrightarrow
 \Pi=N/\overline{\langle\!\langle d_0,d_1\rangle\!\rangle}_N
       \cong\widehat{\pi_1(S_q)}
       \twoheadrightarrow P=\Pi^{(2)}.                    \tag{C.22.2}
\]
They are completion of the finite-cover filling
presentation followed by the maximal pro-two
quotient.  The normal subgroup being killed is
$\alpha$-invariant.  Explicitly, write
$\alpha(c)=f_0^{-1}c^\lambda f_0$ as in C.20.
For $g=1,x$, choose $g'=1$ or $x$ and an integer
$k$ so that
$\rho(\alpha(g)f_0^{-1})=\rho(g'c^k)$.
Then $n=\alpha(g)f_0^{-1}c^{-k}(g')^{-1}\in N$
and
\[
 \alpha(gc^qg^{-1})
       =n\bigl(g'c^q(g')^{-1}\bigr)^\lambda n^{-1}.
\]
The inverse action gives equality of the normal
subgroups.  Thus the actual square induces an
actual automorphism $b$ of $P$ with $b^2=1$.

Pullback of continuous characters identifies
$L=H^1(P,\mathbb Z_2)$ with the submodule of
$V_2=H^1(N,\mathbb Z_2)$ vanishing on $d_0,d_1$.
Indeed this condition is exactly factorization
through both arrows of (C.22.2).  The two cusp
functionals are rotation invariant, so
$W_2=(1-e)V_2$ lies in $L$.  On the deck-trivial
part, transfer identifies this vanishing condition
with $\xi(B)=\xi(C)$ on the basis of $F_0$ from
C.20; multiplication by the odd covering degree
$q$ does not change the condition.  Hence
\[
 L=eL\oplus W_2,\quad\operatorname{rank}eL=2,\quad
 b^*|_{eL}=-I,\quad b^*|_{W_2}=I.                          \tag{C.22.3}
\]
Here the first action follows from $-J^{\mathsf T}$
on $eV_2$ and the equality $\xi(B)=\xi(C)$;
the second is C.21.  Thus $b$ is nontrivial on
integral homology and identity modulo two,
exactly as required in (C.22.1).

The rotation deck action extends geometrically
over both filled cusps and preserves the
unimodular surface cup pairing on $L$.
The averaging idempotent $e$ is self-adjoint:
the adjoint of each deck map is its inverse,
and averaging is unchanged by inversion.
Therefore $eL$ and $W_2$ are orthogonal direct
summands.  The restricted form on each is
unimodular, since the whole form is.  The
nonzero even-rank module $W_2$ consequently
contains integral characters whose reductions
have nonzero cup product.  They are fixed by
$b$ by (C.22.3), contradicting (C.22.1).
This excludes every negative odd coordinate.
C.19 already gives $\lambda_2=1$ and
$\lambda^2=1$, so $\lambda=1$ globally.

Finally, a finite two-subgroup injects on dyadic
homology by C.17.  An element of order $2^a$ with
$a\geq3$ would have a primitive $2^a$th root as
eigenvalue over $\mathbb Q_2$; its minimal
polynomial has degree $2^{a-1}>2$, a contradiction.
Thus all its elements have order one, two or
four.  C.18 and the just-proved quartic assertion
give its stated global multiplier character.
\end{proof}
\statement{C.23}{Proved}{The complete finite-multiplier consequence of LM4 D.42}
Every involution of the full three-cusp loop-outer
group allowing cusp permutations has multiplier
equal to its global dyadic sign.  Together with
C.16 this proves that every finite subgroup has
multiplier image in $\{1,-1\}$, determined by
its dyadic coordinate.  This independently supplies
the transposition inputs LM4 D.36 and D.41.

\begin{proof}
The pure involution case is the refereed R1.
For a transposition, label the fixed cusp $x$ and
the exchanged cusps $y,z=(xy)^{-1}$.
Suppose first that $\lambda_2=1$, and take an
actual involution lift by C.15.  If an odd
coordinate $q$ had $\lambda_q=-1$, the character
\[
 \chi:F\to C_q,\qquad
 \chi(x)=0,\quad\chi(y)=1,\quad\chi(z)=-1
\]
would satisfy $\chi\alpha=\chi$.
Its cover has $q$ degree-one cusps over $x$ and
one degree-$q$ cusp over each of $y,z$.
Schreier generators are
$a_i=y^ixy^{-i}$, $0\leq i<q$, and $b=y^q$;
the last boundary is
$d=z^q=b^{-1}a_{q-1}^{-1}\cdots a_0^{-1}$.
The action permutes the $q$ first cusps by a
translation: writing $\alpha(x)=h^{-1}x^\lambda h$
gives $i\mapsto i-\chi(h)$.  Its order two and
odd $q$ imply $\chi(h)=0$.  Thus all those
cusps are individually fixed.  Filling $b,d$
is the invariant normal quotient
\[
 \ker\chi\longrightarrow
 \widehat{\langle a_0,\ldots,a_{q-1}\mid
                  a_0\cdots a_{q-1}=1\rangle}
       \cong\widehat\pi_{0,q}.
\]
Both maps are verified by killing $b$ and
eliminating $a_{q-1}$; these are completion of
the finite-cover presentation.  Its induced
actual involution is pure with dyadic multiplier
one, hence identity by C.17.  Its nonzero
$q$-adic homology would instead be multiplied
by $-1$.  This contradiction proves the positive
transposition case.

Suppose now $\lambda_2=-1$.  We may choose the
actual lift in the $x$-cusp normalizer.  Indeed
its kernel is $\langle x\rangle=\widehat{\mathbb Z}$
by C.5.  A labelled lift $v$ satisfies
$vxv^{-1}=x^\lambda$, $v^2=x^a$, $(\lambda-1)a=0$.
At negative coordinates $a=0$; at positive
coordinates the prime is odd and division by
two solves the norm correction.  Thus $x^bv$
has order two and gives $\alpha(x)=x^\lambda$.
If there is a positive odd coordinate, choose
the least such prime $q$ and use
\[
 \chi:F\to C_q,\qquad
 \chi(x)=-2,\quad\chi(y)=\chi(z)=1.
\]
Again $\chi\alpha=\chi$.  The cover has genus
$(q-1)/2$ and three cusps.  Filling the two
cusps over $y,z$ gives the exact diagram
\[
\begin{split}
 R&=F/\overline{\langle\!\langle y^q,z^q\rangle\!\rangle}_F
       =C_q\amalg C_q,\\
 K&=\ker(R\to C_q,\ y,z\mapsto1)
       \cong\widehat\pi_{(q-1)/2,1}.
\end{split}
\]
The closed normal subgroup within $\ker\chi$
is the same: $F=(\ker\chi)\langle y\rangle
 =(\ker\chi)\langle z\rangle$, so every
$F$-conjugate of either filled generator is
already a $(\ker\chi)$-conjugate.  The remaining
cusp is $c=x^q$, and the induced actual
involution $\beta$ satisfies $\beta(c)=c^\lambda$.
Conjugation by $y$ gives an actual deck action
$T$ of order $q$, whose homology commutes
with $A=\beta_*$.  Schreier generators
$u_i=y^izy^{-(i+1)}$ have the sole relation
$u_0\cdots u_{q-1}=1$ and satisfy $T(u_i)=u_{i+1}$.
Hence for every $p$,
\[
 H_1(K,\mathbb Q_p)\cong\mathbb Q_p[X]/\Phi_q(X),
            \qquad A^2=1,\quad AT=TA.
\]
The one-boundary commutator bivector gives a
nondegenerate alternating form with multiplier
$\lambda_p$ under $A$ and multiplier one under $T$.

If $p\ne q$ has even order $d$ modulo $q$,
each irreducible Frobenius orbit of primitive
$q$th roots over $\mathbb Q_p$ contains their
inverses, since $p^{d/2}=-1\bmod q$.
Every such irreducible summand occurs once and
is nondegenerate for the alternating form.
The commutant on it is a field, so $A^2=1$
makes $A$ the scalar $1$ or $-1$ there.
Thus $\lambda_p=1$.
There is a prime divisor $p$ of $q-1$ with
even order modulo $q$: otherwise every prime
factor of $q-1$ belongs to the odd-order
subgroup of $(\mathbb Z/q)^\times$, while their
product with multiplicities is $q-1=-1$ of
order two.  Such a $p$ is smaller than $q$.
For $p=2$ its sign is negative by hypothesis;
for odd $p<q$ it is negative by minimality.
Both contradict $\lambda_p=1$.  This proves
the negative transposition assertion.

For a finite subgroup, each element's image
in $\mathfrak S_3$ has order $m=1,2$ or $3$,
and its $m$th power belongs to the pure group.
R1 makes its order divide $2m$.  Order four
is impossible: its square would be a pure
involution of multiplier $-1$, requiring an
odd dyadic unit whose square is $-1$, whereas
every odd dyadic square is $1\bmod8$.
The remaining orders are $1,2,3,6$.
Order three has multiplier one by C.16.
For order six, its square has multiplier one
and its cube has global multiplier $1$ or
$-1$, so its own multiplier equals that of
its cube.  This proves the full consequence.
\end{proof}

